\documentclass[11pt]{article}

\usepackage[margin=1in]{geometry}
\usepackage{amsmath, amssymb}
\usepackage[T1]{fontenc}
\usepackage{lmodern}
\usepackage{microtype}
\usepackage{graphicx}
\usepackage{booktabs}
\usepackage{array}
\usepackage{float}
\usepackage{xcolor}
\usepackage{fancyhdr}
\usepackage[font=small,labelfont=bf,skip=5pt]{caption}
\usepackage{listings}
\lstdefinestyle{pycode}{
  language=Python, basicstyle=\ttfamily\scriptsize,
  keywordstyle=\color{psink}\bfseries, commentstyle=\color{psaccent}\itshape,
  stringstyle=\color{psink}, numbers=left, numberstyle=\tiny\color{psrule},
  numbersep=7pt, breaklines=true, breakatwhitespace=true, showstringspaces=false,
  frame=none, xleftmargin=20pt, aboveskip=6pt, belowskip=6pt, columns=fullflexible,
  keepspaces=true, upquote=true,
}
\lstset{style=pycode}

\definecolor{psink}{RGB}{23,54,93}
\definecolor{psaccent}{RGB}{192,80,77}
\definecolor{psrule}{RGB}{200,209,222}

\usepackage[colorlinks=true,allcolors=psink]{hyperref}

\graphicspath{{./}}
\emergencystretch=3em

\setlength{\textfloatsep}{10pt plus 2pt minus 2pt}
\setlength{\intextsep}{10pt plus 2pt minus 2pt}
\setlength{\abovecaptionskip}{4pt}
\setlength{\belowcaptionskip}{0pt}

\makeatletter
\renewcommand\section{\@startsection{section}{1}{0pt}%
  {-2.6ex \@plus -1ex \@minus -.2ex}{1.3ex \@plus .2ex}%
  {\normalfont\scshape\large\color{psink}}}   % lmr has no bold small caps
\renewcommand\subsection{\@startsection{subsection}{2}{0pt}%
  {-2.2ex \@plus -1ex \@minus -.2ex}{1.0ex \@plus .2ex}%
  {\normalfont\bfseries\normalsize\color{psink}}}
\renewcommand\subsubsection{\@startsection{subsubsection}{3}{0pt}%
  {-1.9ex \@plus -.8ex \@minus -.2ex}{0.8ex \@plus .2ex}%
  {\normalfont\bfseries\itshape\normalsize\color{psink}}}
\renewcommand\subparagraph{\@startsection{subparagraph}{5}{0pt}%
  {-1.6ex \@plus -.6ex \@minus -.2ex}{0.5ex \@plus .1ex}%
  {\normalfont\itshape\normalsize\color{psink}}}
\renewcommand\paragraph{\@startsection{paragraph}{4}{0pt}%
  {1.7ex \@plus 1ex \@minus .2ex}{-0.9em}%
  {\normalfont\bfseries\normalsize\color{psink}}}
\makeatother

\newenvironment{tight}
  {\begin{itemize}\setlength{\itemsep}{3pt}\setlength{\parskip}{0pt}%
   \setlength{\topsep}{3pt}\setlength{\leftmargini}{1.4em}}
  {\end{itemize}}

\setcounter{secnumdepth}{0}   % headings carry the assignment's own labels
\setlength{\headheight}{14pt}
\pagestyle{fancy}
\fancyhf{}
\fancyhead[L]{\footnotesize\scshape\color{psink} MGT 595 \textemdash\ Problem Set 5}
\fancyhead[R]{\footnotesize\color{psink}\thepage}
\renewcommand{\headrulewidth}{0.4pt}
\renewcommand{\headrule}{\hbox to\headwidth{%
  \color{psrule}\leaders\hrule height \headrulewidth\hfill}}

\begin{document}

\thispagestyle{empty}
\begin{flushleft}
{\bfseries\LARGE\color{psink} Quantitative Investing}\\[3pt]
{\scshape\Large\color{psink} Problem Set 5}\\[10pt]
{\color{psrule}\rule{\textwidth}{1.2pt}}\\[5pt]
{\small Ashley Wen\,(xw564)\quad\textbar\quad Emily Wu\,(esw63)%
\quad\textbar\quad Grace Gu\,(xg325)\hfill \today}
\end{flushleft}

\vspace{1.4em}

{\color{psink}\textbf{Contents.}} Part I studies how the factor and corner-portfolio
returns vary over the business cycle (Questions (a)--(b)). Part II runs the
characteristics-versus-covariances horse race, first on the 25 size and BE/ME
portfolios (Questions (c)--(e)) and then on the 25 size and momentum portfolios
(Questions (f)--(h)).

\vspace{0.8em}

% =====================================================================
\section{Data and Conventions}
% =====================================================================

\begin{tight}
  \item \textbf{Sample.} Monthly data from 1926-07 to 2026-06 (UMD and the momentum
  portfolios start 1927-01). The NBER dummy marks 201 recession months across 16 recessions.
  \item \textbf{Returns.} All figures are percent per month. Portfolio returns are in
  excess of the T-bill; RMRF is already excess, and SMB, HML and UMD are zero-cost
  long--short returns, so they are used as given.
  \item \textbf{Inference.} Residuals cluster in recession episodes, so the $t$-statistics in
  Part I use Newey--West standard errors with 12 lags. Fama--MacBeth standard errors in
  Part II are the time-series standard deviation of the monthly $\hat\gamma_t$ divided
  by $\sqrt{T}$.
  \item \textbf{Shared code.} All parts load data through \texttt{common/data.py} and run
  Fama--MacBeth through \texttt{common/fmb.py}, so lags, samples and beta definitions
  are identical across Questions (a)--(h).
\end{tight}

% =====================================================================
\section{Part I: Business Cycle Variation}
% =====================================================================

\subsection{Question (a)}

We estimate $R_t = a + b\,D_t + \varepsilon_t$, where $D_t = 1$ in NBER recession months,
so $a$ is the non-recession mean and $a+b$ the recession mean.

\begin{table}[H]\centering\small
\caption{Factor returns in and out of recessions (\% per month)}
\begin{tabular}{lcccc}
\toprule
 & Non-recession $a$ & Recession $a+b$ & Difference $b$ & $t(b)$ \\
\midrule
RMRF & 0.93 & $-0.47$ & $-1.40$ & $-2.50$ \\
SMB  & 0.19 & 0.03    & $-0.16$ & $-0.82$ \\
HML  & 0.40 & 0.15    & $-0.25$ & $-0.75$ \\
UMD  & 0.71 & 0.22    & $-0.49$ & $-1.00$ \\
\bottomrule
\end{tabular}
\par\smallskip\footnotesize Newey--West (12 lags) $t$-statistics. UMD starts 1927-01.
\end{table}

\begin{tight}
  \item \textbf{What does this tell you?} Only the market is reliably worse in
  recessions; SMB, HML and UMD give up some return but none of the drops is
  distinguishable from zero. All three factors still earn a positive average return in
  recession months, while the market loses money. So the factor premia are not
  mainly a reward for losing money in recessions, at least not in the way the market premium is.
  \item \textbf{Based on these results, do you believe SMB or HML is more consistent
  with a risk-based explanation?} SMB is, but only weakly. Its premium all but
  disappears in recessions, so small stocks fail to pay off exactly when times are
  bad, which is what a risk story needs. HML instead keeps a positive return in
  recessions, in line with Lakonishok, Shleifer and Vishny's (1994) finding that value does not
  underperform in bad states. With 201 recession months neither difference is
  significant, so we read this as a lean rather than a verdict.
\end{tight}

\subsection{Question (b)}

\begin{table}[H]\centering\small
\caption{Corner portfolio excess returns in and out of recessions (\% per month)}
\begin{tabular}{lcccc}
\toprule
 & Non-recession $a$ & Recession $a+b$ & Difference $b$ & $t(b)$ \\
\midrule
\multicolumn{5}{l}{\textit{Size--BE/ME corners}} \\
Small Growth & 0.83 & $-0.63$ & $-1.46$ & $-1.30$ \\
Small Value  & 1.69 & $-0.38$ & $-2.07$ & $-2.62$ \\
Big Growth   & 0.90 & $-0.34$ & $-1.25$ & $-2.21$ \\
Big Value    & 1.24 & $-0.12$ & $-1.35$ & $-1.51$ \\
\addlinespace
\multicolumn{5}{l}{\textit{Size--momentum corners}} \\
Small Losers  & 0.80 & $-0.26$ & $-1.06$ & $-1.17$ \\
Small Winners & 1.84 & $-0.22$ & $-2.06$ & $-2.98$ \\
Big Losers    & 0.58 & $-0.83$ & $-1.42$ & $-1.66$ \\
Big Winners   & 1.24 & $-0.40$ & $-1.64$ & $-2.73$ \\
\addlinespace
\multicolumn{5}{l}{\textit{Within-size spreads}} \\
Small Value $-$ Growth   & 0.87 & 0.25 & $-0.62$ & $-1.03$ \\
Big Value $-$ Growth     & 0.34 & 0.23 & $-0.11$ & $-0.14$ \\
Small Winners $-$ Losers & 1.04 & 0.04 & $-1.00$ & $-1.29$ \\
Big Winners $-$ Losers   & 0.64 & 0.39 & $-0.25$ & $-0.40$ \\
\bottomrule
\end{tabular}
\par\smallskip\footnotesize Excess returns over the T-bill. Newey--West (12 lags)
$t$-statistics. Momentum portfolios start 1927-01.
\end{table}

\begin{tight}
  \item \textbf{Do you find stronger cyclical patterns for these extreme portfolios
  than for the factors?} Yes, but mostly because the corners are long-only portfolios that
  carry the market's recession losses. Every corner falls by one to two percent per
  month in recessions, and the high-return corners (Small Value and both
  Winner portfolios) fall the most. Once we net out the market by taking
  spreads within a size group, the pattern shrinks back to factor-sized, insignificant
  drops that stay positive in recessions. The extreme portfolios are more cyclical,
  but the value and momentum premia themselves are not.
\end{tight}

% =====================================================================
\section{Part II: Characteristics vs.\ Covariances}
% =====================================================================

\subsection{25 Size and BE/ME Portfolios}

\subsubsection{Question (c)}

\paragraph{How the PS3 code was extended.} We used an LLM to produce most of the PS5
code extension, starting from our PS3 implementation in cells 1(b) and 1(f). The main
changes were adding joint regressions on RMRF, SMB and HML to estimate the additional
factor betas, generalizing the monthly cross-sectional regression into a function that
runs all three specifications, and producing a side-by-side results table. The extension
preserved the original FMB standard-error calculation, fixed betas across months, and
lagged, time-varying characteristics. The final implementation passes the included
verification checks and reproduces the PS3 size and BE/ME estimates to four decimal
places when using the original sample and market-only beta. Equation (1) retains the PS3
market-only beta, while equations (2) and (3) use joint three-factor betas, following our
stated team convention.

\paragraph{Were the changes correct?} Yes. All checks in Cell 7 of our notebook pass:
\begin{tight}
  \item Equation (2) equals a single regression of average returns on the fixed betas.
  \item A numpy re-run of one month confirms which beta each equation uses.
  \item All equations use the same 1,194 months with 25 portfolios each.
  \item Every SE equals $sd(\hat\gamma_t)/\sqrt{T}$.
\end{tight}

\paragraph{Fixed betas, time-varying characteristics.} The betas are estimated once and
stay fixed across all months of the second pass. Size and BE/ME are lagged and change
over time. Checks 1--3 in Cell 7 confirm this structure was preserved.

\begin{table}[H]\centering\small
\caption{Fama--MacBeth estimates, 25 Size--BE/ME portfolios, full sample (\% per month)}
\begin{tabular}{lccc}
\toprule
 & (1) Characteristics & (2) Covariances & (3) Combined \\
\midrule
$\gamma_0$       & 1.7547 (4.15)    & 1.6433 (4.16)    & 1.4108 (3.53)    \\
$\gamma_M$       & $-0.4777$ ($-1.34$) & $-0.9042$ ($-2.21$) & 0.2049 (0.52) \\
$\gamma_{size}$  & $-0.0674$ ($-2.30$) & --            & $-0.0932$ ($-1.96$) \\
$\gamma_{B/M}$   & 0.2327 (4.08)    & --               & 0.4808 (4.39)    \\
$\gamma_{SMB}$   & --               & 0.1087 (1.12)    & $-0.1567$ ($-0.83$) \\
$\gamma_{HML}$   & --               & 0.3752 (3.54)    & $-0.4656$ ($-2.05$) \\
\midrule
Avg.\ $R^2$      & 0.490            & 0.485            & 0.573            \\
\bottomrule
\end{tabular}
\par\smallskip\footnotesize Average monthly FMB coefficients with $t$-statistics (testing
$\gamma=0$) in parentheses. 1927-01 to 2026-06, $T=1194$, $N=25$. $\beta_M$ in equation
(1) is the market-only (CAPM) beta; in equations (2) and (3) it is the joint FF3 loading,
so the two are different numbers. $\ln(size)$ is lagged one month; $\ln(BE/ME)$ is the
previous calendar year's value.
\end{table}

\paragraph{Check against PS3.} On 1969-07 to 2026-06 ($T = 684$), with the market-only
beta re-estimated on that sample, equation (1) gives $\gamma_{size} = -0.0194$
($t = -0.64$) and $\gamma_{B/M} = 0.1573$ ($t = 2.11$). These match our PS3 results
exactly to 4 d.p., and $\gamma_0$ and $\gamma_M$ also match. The full-sample equation
(1) estimates differ because the sample period is different.

\paragraph{Implementation qualifications.}
\begin{tight}
  \item \emph{Market beta in equation (1):} the wording is ambiguous (see the sensitivity
  check below). With $\beta_M^{\mathrm{FF3}}$, $\gamma_{size}$ changes from $-0.0674$ ($t = -2.30$)
  to $-0.0367$ ($t = -1.20$), while $\gamma_{B/M}$ stays at $t > 3$.
  \item \emph{BE/ME timing:} this follows the PS3/team calendar-year rule, not the
  Fama--French July--June timing described in Lecture 2.
  \item \emph{Generated regressors:} the betas are estimated on the full sample, so they
  contain measurement error and a look-ahead element. The FMB SEs are not adjusted for
  beta-estimation error.
  \item \emph{PS3 comparison precision:} the match is confirmed to 4 d.p., which is all
  PS3 printed.
\end{tight}

\paragraph{Sensitivity check: which market beta to use in equation (1).}
Step 1 of Question (c) says to estimate $\beta_{iM}$, $\beta_{iSMB}$ and $\beta_{iHML}$
``using full-period OLS regressions on RMRF, SMB, and HML,'' and the same symbol
$\beta_{iM}$ appears in all three equations. This can be read as one joint FF3 regression
whose market loading is used everywhere. However, the same step notes that ``$\beta_{iM}$
was already estimated in Problem Set 3,'' and we are asked to verify that equation (1)
matches the PS3 $\gamma_{size}$ and $\gamma_{B/M}$. In PS3, $\beta_{iM}$ came from a
market-only regression. The assignment does not settle which reading is intended. The
team convention uses the market-only beta $\beta_M^{\mathrm{CAPM}}$ in (1) and the joint FF3
betas ($\beta_M^{\mathrm{FF3}}$, $\beta_{SMB}$, $\beta_{HML}$) in (2) and (3). Under the PS3 setup
(Cell 7), this reproduces the PS3 estimates and $t$-statistics to 4 d.p. Table~\ref{tab:sens}
shows what changes with $\beta_M^{\mathrm{FF3}}$ in equation (1); the sample (1927-01 to 2026-06,
$T = 1194$), the characteristics and the FMB standard-error formula are unchanged.

\begin{table}[H]\centering\small
\caption{Equation (1) under the two market-beta definitions (\% per month)}\label{tab:sens}
\begin{tabular}{lcc}
\toprule
 & $\beta_M^{\mathrm{CAPM}}$ (main) & $\beta_M^{\mathrm{FF3}}$ (sensitivity) \\
\midrule
$\gamma_0$      & 1.7547 (4.15)       & 1.5272 (3.76)       \\
$\gamma_M$      & $-0.4777$ ($-1.34$) & $-0.5264$ ($-1.21$) \\
$\gamma_{size}$ & $-0.0674$ ($-2.30$) & $-0.0367$ ($-1.20$) \\
$\gamma_{B/M}$  & 0.2327 (4.08)       & 0.2020 (3.57)       \\
\bottomrule
\end{tabular}
\par\smallskip\footnotesize $t$-statistics in parentheses. 1927-01 to 2026-06, $T=1194$.
\end{table}

Average monthly $R^2$ moves from 0.4898 to 0.4852, and the two market betas have a
cross-sectional correlation of 0.65. The main difference is $\gamma_{size}$: its $|t|$ is
above 1.96 with $\beta_M^{\mathrm{CAPM}}$ but not with $\beta_M^{\mathrm{FF3}}$. $\gamma_{B/M}$ has $t$
above 3 in both versions.

\subsubsection{Question (d)}

\paragraph{Step 1 --- Preliminary interpretation.} Equation (3) has the highest average
monthly cross-sectional $R^2$ (0.573, versus 0.490 for equation (1) and 0.485 for
equation (2)), although this measures in-sample fit and partly reflects the additional
regressors, so it does not by itself distinguish between the characteristics and risk
explanations. In the combined specification, BE/ME stays strongly significant, and its
coefficient roughly doubles relative to equation (1) (0.4808, $t = 4.39$, versus 0.2327,
$t = 4.08$). Size remains negative but narrowly misses the 5\% level ($t = -1.96$,
$p = 0.0508$). By contrast, $\gamma_{HML}$ turns from positive and significant in
equation (2) (0.3752, $t = 3.54$) to negative in equation (3) ($-0.4656$, $t = -2.05$),
and $\gamma_{SMB}$ is insignificant in both equations ($t = -0.83$ in (3)). Taken
together, BE/ME keeps explanatory power after controlling for factor exposures, while the
betas lose their expected positive sign, which leans toward a characteristics
interpretation; however, the roughly doubled standard errors in equation (3) suggest
betas and characteristics are hard to separate, so we treat this conclusion as tentative.

\paragraph{Step 2 --- LLM discussion.}
\begin{tight}
  \item \textbf{``What is the strongest economic argument that size and BE/ME are priced
  because they are characteristics rather than because they co-vary with systematic
  risk?''} The strongest argument is that investors may systematically misprice firms
  with certain observable characteristics. For example, investors may overvalue glamorous,
  low-BE/ME firms and undervalue distressed or neglected high-BE/ME firms, so the later
  return premium reflects price correction rather than compensation for bearing systematic
  risk. Under this view, size and BE/ME themselves should predict returns even after
  controlling for factor betas.
  \item \textbf{``What is the strongest economic argument that size and BE/ME premia are
  compensation for systematic risk --- and what would that imply about the betas on SMB
  and HML?''} The risk view says small and high-BE/ME firms tend to perform poorly in
  certain bad economic states, so investors require higher expected returns to hold them.
  SMB and HML are then factor-mimicking portfolios capturing these common systematic
  risks. Therefore, firms with higher $\beta_{SMB}$ or $\beta_{HML}$ should earn higher
  average returns, and once those betas are controlled for, size and BE/ME should have
  little additional explanatory power.
  \item \textbf{``In a Fama--MacBeth horse race between equations (1), (2), and (3), what
  empirical pattern would cleanly support the characteristics view? What pattern would
  support the risk view?''} Suppose equation (1) contains characteristics, equation (2)
  contains factor betas, and equation (3) includes both:
  \begin{tight}
    \item Characteristics view: size and BE/ME are significant in (1) and remain
    significant in (3), while $\gamma_{SMB}$ and $\gamma_{HML}$ become weak or insignificant.
    \item Risk view: $\gamma_{SMB}$ and $\gamma_{HML}$ are significant in (2) and remain
    significant in (3), while $\gamma_{size}$ and $\gamma_{B/M}$ become weak or insignificant.
    \item If both remain significant in (3), the evidence is mixed rather than a clean
    victory for either explanation.
  \end{tight}
\end{tight}

\paragraph{Step 3 --- Evaluation.}
\begin{tight}
  \item \textbf{Does the LLM correctly identify the key distinguishing prediction --- that
  if characteristics dominate in equation (3), it is hard to sustain a pure risk story?}
  The LLM correctly identifies that characteristics remaining significant in equation (3),
  while beta premia weaken, would challenge a pure SMB/HML risk explanation. However,
  significance alone is not enough and the signs of the coefficients also matter.
  \item \textbf{Does the LLM correctly note that the horse race is inherently noisy because
  characteristics and betas are correlated by construction (size sorts produce size beta
  differences)?} The LLM's answer does not mention that characteristics and betas are
  strongly correlated because of the portfolio sorts. Our notebook shows correlations of
  $-0.974$ between average log size and $\beta_{SMB}$ and 0.948 between average log BE/ME
  and $\beta_{HML}$, making their separate effects difficult to identify.
  \item \textbf{Does your empirical evidence from equations (1)--(3) more strongly support
  one side? Be explicit about which coefficients drive your conclusion.} Our full-sample
  results lean toward the characteristics view, especially for BE/ME: $\gamma_{B/M}$
  increases from 0.2327 ($t=4.08$) in equation (1) to 0.4808 ($t=4.39$) in equation (3).
  $\gamma_{size}$ remains negative but narrowly misses 5\% significance ($p=0.0508$), while
  $\gamma_{SMB}$ is insignificant in both equations (2) and (3). $\gamma_{HML}$ remains
  significant but reverses from positive to negative ($-0.4656$, $t=-2.05$), which
  conflicts with the positive premium predicted by the LLM's risk explanation. The evidence
  therefore favors characteristics, but the strong correlations make this conclusion
  tentative.
\end{tight}

\subsubsection{Question (e)}

\begin{table}[H]\centering\small
\caption{Fama--MacBeth estimates, 25 Size--BE/ME portfolios, 1963-01 onward (\% per month)}
\begin{tabular}{lccc}
\toprule
 & (1) Characteristics & (2) Covariances & (3) Combined \\
\midrule
$\gamma_0$       & 1.2790 (3.97)       & 1.1429 (4.42)       & 1.3700 (3.58)       \\
$\gamma_M$       & $-0.2247$ ($-0.77$) & $-0.5261$ ($-1.73$) & $-0.2827$ ($-0.78$) \\
$\gamma_{size}$  & $-0.0437$ ($-1.52$) & --                  & 0.0045 (0.09)       \\
$\gamma_{B/M}$   & 0.1704 (2.50)       & --                  & 0.4876 (3.57)       \\
$\gamma_{SMB}$   & --                  & 0.1228 (1.10)       & 0.0527 (0.26)       \\
$\gamma_{HML}$   & --                  & 0.3364 (3.06)       & $-0.4949$ ($-1.87$) \\
\midrule
Avg.\ $R^2$      & 0.525               & 0.528               & 0.593               \\
\bottomrule
\end{tabular}
\par\smallskip\footnotesize $t$-statistics in parentheses. 1963-01 to 2026-06, $T = 762$,
$N = 25$, with betas re-estimated on the same months (team conventions). $\beta_M$ is the
CAPM beta in (1) and the joint FF3 loading in (2) and (3).
\end{table}

\begin{table}[H]\centering\small
\caption{Full sample versus post-1963, 25 Size--BE/ME portfolios}
\begin{tabular}{lccc}
\toprule
 & Equation & Full: Est.\ ($t$) & Post-1963: Est.\ ($t$) \\
\midrule
$\gamma_{size}$ & (1) & $-0.0674$ ($-2.30$) & $-0.0437$ ($-1.52$) \\
$\gamma_{B/M}$  & (1) & 0.2327 (4.08)       & 0.1704 (2.50)       \\
$\gamma_{SMB}$  & (2) & 0.1087 (1.12)       & 0.1228 (1.10)       \\
$\gamma_{HML}$  & (2) & 0.3752 (3.54)       & 0.3364 (3.06)       \\
$\gamma_{size}$ & (3) & $-0.0932$ ($-1.96$) & 0.0045 (0.09)       \\
$\gamma_{B/M}$  & (3) & 0.4808 (4.39)       & 0.4876 (3.57)       \\
$\gamma_{SMB}$  & (3) & $-0.1567$ ($-0.83$) & 0.0527 (0.26)       \\
$\gamma_{HML}$  & (3) & $-0.4656$ ($-2.05$) & $-0.4949$ ($-1.87$) \\
\bottomrule
\end{tabular}
\par\smallskip\footnotesize Full sample: 1927-01 to 2026-06 ($T = 1194$). Post-1963:
1963-01 to 2026-06 ($T = 762$).
\end{table}

\begin{tight}
  \item \textbf{Does the characteristics-vs.-covariances result change for the post-1963
  sample?} For BE/ME, the result is basically the same. In equation (3), $\gamma_{B/M}$ is
  almost unchanged (0.4876, $t = 3.57$), while $\gamma_{HML}$ stays negative and is now
  insignificant at 5\% ($t = -1.87$, $p = 0.06$). So the value evidence still leans toward
  the characteristics view. The size result changes: $\gamma_{size}$ is no longer
  significant in equation (1) ($t = -1.52$) and is essentially zero in equation (3)
  ($t = 0.09$), and $\gamma_{SMB}$ is insignificant in both samples. After 1963, neither
  size nor its beta has a statistically significant coefficient in the combined
  specification, so the size dimension gives no evidence for either side.
  \item \textbf{Do you feel more or less confident in your conclusion from part (d)?} We
  are more confident about BE/ME, because $\gamma_{B/M}$ in equation (3) has a similar size
  and stays significant in both samples. We are less confident about size, which was
  already borderline in part (d) and disappears after 1963. Overall the conclusion is
  still tentative. With fewer months, the standard errors in equation (3) are larger
  (e.g.\ the SE of $\gamma_{HML}$ is 0.2644 vs.\ 0.2270), and betas and characteristics
  are still strongly correlated by construction of the 25 portfolios.
  \item \textbf{Why might the sample period matter for this particular question?}
  \begin{tight}
    \item \emph{Data quality:} before 1963, book equity is not from Compustat and covers
    fewer firms, so BE/ME and the portfolios may be measured less precisely.
    \item \emph{In-sample vs.\ out-of-sample:} the post-1963 period is close to the sample
    in which the size and value effects were originally documented, so the pre-1963 data
    works partly as an out-of-sample check.
    \item \emph{Changing premia and risk:} the size premium seems to depend on the period.
    The pre-1963 sample also includes the Great Depression, when distress risk was
    especially important. If factor loadings change over time, full-sample fixed betas fit
    some periods worse than others.
  \end{tight}
\end{tight}

\subsection{25 Size and Momentum Portfolios}

\subsubsection{Question (f)}

\paragraph{Method.} We estimate three Fama--MacBeth (FMB) specifications using the 25
size- and momentum-sorted portfolios. Equation (1) includes market beta, log size, and
past returns (ret212); equation (2) includes market, SMB, and UMD betas; equation (3)
combines characteristics and factor betas. First-stage betas are estimated from
time-series regressions of portfolio excess returns on the relevant factors. Following
the shared convention, equation (1) uses the CAPM market beta estimated from RMRF alone,
while equations (2) and (3) use market, SMB, and UMD betas estimated jointly. Betas
remain fixed in the second stage, while $\ln(size)$ and ret212 are lagged one month and
vary over time. Each month, we run the relevant cross-sectional regression and then
average the monthly coefficients. FMB standard errors equal the time-series standard
deviation of the monthly coefficients divided by the square root of the number of months.
All equations use the same 1,186 months from February 1927 to June 2026.

\begin{table}[H]\centering\small
\caption{Fama--MacBeth estimates, 25 Size--momentum portfolios, full sample (\% per month)}\label{tab:momfull}
\begin{tabular}{lccc}
\toprule
 & (1) Characteristics & (2) Covariances & (3) Combined \\
\midrule
$\gamma_0$        & 2.0220 (5.51)       & 1.5006 (4.55)       & 2.2296 (4.91)       \\
$\gamma_M$        & $-0.5967$ ($-2.21$) & $-0.6567$ ($-1.89$) & $-0.9362$ ($-2.54$) \\
$\gamma_{size}$   & $-0.1319$ ($-4.35$) & --                  & $-0.0670$ ($-1.53$) \\
$\gamma_{ret212}$ & 0.0068 (4.49)       & --                  & 0.0076 (2.77)       \\
$\gamma_{SMB}$    & --                  & 0.3616 (3.45)       & 0.1031 (0.56)       \\
$\gamma_{UMD}$    & --                  & 0.6179 (4.47)       & 0.0745 (0.35)       \\
\midrule
Avg.\ $R^2$       & 0.531               & 0.543               & 0.620               \\
\bottomrule
\end{tabular}
\par\smallskip\footnotesize Average monthly FMB coefficients with $t$-statistics in
parentheses. Excess returns in percent per month; 1927-02 to 2026-06, $T = 1186$.
$\beta_M$ is the CAPM beta in (1) and the joint RMRF/SMB/UMD loading in (2) and (3).
ret212 is the portfolio's past 2--12 month return in percent, lagged one month, so
$\gamma_{ret212}$ is per percentage point of past return: a portfolio whose ret212 is
10 percentage points higher is predicted to earn about 0.07\% more per month.
\end{table}

\paragraph{Interpretation.} In equation (1), both smaller size and higher past returns
predict higher subsequent excess returns. In equation (2), $\gamma_{SMB}$ and
$\gamma_{UMD}$ are significant when characteristics are excluded. However, in the
combined equation (3), $\gamma_{ret212}$ remains significant ($t = 2.77$), while
$\gamma_{SMB}$ ($t = 0.56$), $\gamma_{UMD}$ ($t = 0.35$), and $\gamma_{size}$
($t = -1.53$) are not. This favors the characteristics interpretation for momentum over a
pure explanation based on the measured SMB and UMD exposures, although collinearity and
first-stage beta estimation error limit how decisively the two explanations can be
separated.

\subsubsection{Question (g)}

\paragraph{Step 1 --- Preliminary interpretation.} Before consulting an LLM, our
interpretation was that both explanations had some support when tested separately, but
the combined model favored momentum characteristics. Size and past returns are
significant in equation (1), while $\gamma_{SMB}$ and $\gamma_{UMD}$ are significant in
equation (2). Once they compete in equation (3), $\gamma_{ret212}$ remains significant,
while $\gamma_{SMB}$ and $\gamma_{UMD}$ lose significance. This leans toward the
characteristics view, although correlated characteristics and factor exposures prevent a
definitive causal conclusion. \emph{We recorded the qualitative view above before the LLM
discussion; the numerical values were later updated to match the group conventions.}

\paragraph{Step 2 --- LLM interaction.} We asked: ``For the momentum premium, is the
debate between a characteristics explanation and a risk explanation the same as for the
value premium? What makes momentum harder to reconcile with conventional risk-based
models than value?''

The LLM argued that both value and momentum raise the same broad question of systematic
risk versus mispricing. Value has potential links to persistent fundamental risks such as
distress, while momentum is based on recent price performance and lacks an equally
obvious stable economic risk exposure. It also noted that conventional models have
difficulty explaining momentum's return continuation, while behavioral mechanisms such as
investor underreaction provide a natural alternative. Momentum also has severe
state-dependent crash risk, particularly when past losers rebound sharply after prolonged
declines. It also correctly stated that the existence of crash risk alone does not prove
that momentum's average premium is compensation for systematic risk.

\paragraph{Step 3 --- Evaluation.}
\begin{tight}
  \item \textbf{Does the LLM recognize that momentum has severe, state-dependent crash
  risk, and correctly explain why the existence of crash risk alone does not establish
  that momentum's average return is risk compensation?} Partly. It recognizes momentum's
  state-dependent crashes and correctly avoids the shortcut that a risky strategy must
  earn its premium because of that risk. However, it does not fully explain why crashes
  alone are insufficient evidence for a risk premium: risk compensation requires losses to
  occur systematically in economically bad states, whereas momentum crashes often occur
  during sharp market rebounds.
  \item \textbf{Does the LLM mention behavioral explanations for momentum, and correctly
  contrast these with the risk-based story for value?} Yes, though incompletely. It
  mentions underreaction but not other behavioral mechanisms such as overreaction or
  investor inattention. Importantly, it does not simply import the value risk story into
  momentum: it contrasts value's potential links to persistent fundamental risks such as
  distress with momentum, which lacks an equally obvious stable fundamental-risk exposure.
  \item \textbf{Does your empirical result for momentum differ from your value result in
  part (d)? Does the LLM's framing help explain the difference?} Our results provide
  evidence for both sides. The strongest evidence for the risk view appears in equation
  (2): $\gamma_{UMD}$ is 0.618\% per month, very close to UMD's average return of about
  0.60\% per month over the same 1,186 months, which is consistent with a risk
  interpretation when characteristics are excluded. Momentum crashes and time-varying
  exposures also leave room for conditional risk explanations. However, the
  characteristics evidence becomes stronger in equation (3). $\gamma_{ret212}$ remains
  significant ($t = 2.77$), while $\gamma_{UMD}$ falls from 0.618\% to 0.075\% per month. The equation
  (3) UMD risk price is far below UMD's average return of 0.60\% per month and is itself
  insignificant ($t = 0.35$), suggesting that UMD beta no longer commands the return
  predicted by the simple traded-factor risk story once past-return characteristics are
  included. $\gamma_{SMB}$
  is also insignificant. The beta-only model is not a complete successful risk model
  either: its intercept remains large and significant, while $\gamma_M$ is negative.
  Overall, the evidence leans toward characteristics, but it does not prove that momentum
  is purely behavioral because characteristics and factor exposures are highly correlated
  and more complicated risk models could still matter.

  The value result is similar but not identical. In both cases, the
  characteristic retains explanatory power in the combined specification while the
  corresponding factor-beta evidence becomes less supportive of a conventional positive
  risk-premium story. For value, $\gamma_{B/M}$ remains strongly significant in equation
  (3) (0.4808, $t = 4.39$), while $\gamma_{SMB}$ is insignificant ($t = -0.83$).
  $\gamma_{HML}$ remains significant but becomes negative ($-0.4656$, $t = -2.05$), which
  is inconsistent with the positive premium predicted by the standard HML risk story. For
  momentum, $\gamma_{ret212}$ remains significant ($t = 2.77$), while $\gamma_{SMB}$
  ($t = 0.56$) and $\gamma_{UMD}$ ($t = 0.35$) are insignificant, although $\gamma_M$ is
  significantly negative ($t = -2.54$). Thus, neither model has a perfectly clean
  risk-pricing pattern, and momentum does not provide a stronger or better-identified
  characteristics result. The correlation between $\beta_{UMD}$ and the
  momentum characteristic is very high, just as $\beta_{HML}$ and BE/ME are highly
  correlated for value, and the full-sample $\gamma_{ret212}$ $t$-statistic of 2.77 is
  lower than the 4.39 of $\gamma_{B/M}$.

  The LLM's framing helps explain the difference only partly. Momentum has less obvious
  fundamental-risk foundations than value, making its premium harder to reconcile with
  conventional risk models, so it is less surprising that UMD beta is not priced once the
  past-return characteristic is included.
\end{tight}

\subsubsection{Question (h)}

\paragraph{Method.} We repeat the three specifications using return months from January
1963 onward and re-estimate all first-stage betas within the restricted sample. Lagged
characteristics are constructed before the sample is restricted, so January 1963 uses
information available in December 1962. The restricted sample contains 762 months from
January 1963 to June 2026.

\begin{table}[H]\centering\small
\caption{Fama--MacBeth estimates, 25 Size--momentum portfolios, 1963-01 onward (\% per month)}
\begin{tabular}{lccc}
\toprule
 & (1) Characteristics & (2) Covariances & (3) Combined \\
\midrule
$\gamma_0$        & 1.4933 (5.11)       & 1.1499 (4.83)       & 1.2531 (3.82)       \\
$\gamma_M$        & $-0.5506$ ($-2.18$) & $-0.4002$ ($-1.43$) & $-0.9542$ ($-2.55$) \\
$\gamma_{size}$   & $-0.0578$ ($-2.20$) & --                  & 0.0102 (0.19)       \\
$\gamma_{ret212}$ & 0.0075 (4.98)       & --                  & 0.0080 (4.13)       \\
$\gamma_{SMB}$    & --                  & 0.1734 (1.47)       & 0.1577 (0.68)       \\
$\gamma_{UMD}$    & --                  & 0.6659 (4.35)       & $-0.0427$ ($-0.20$) \\
\midrule
Avg.\ $R^2$       & 0.574               & 0.591               & 0.658               \\
\bottomrule
\end{tabular}
\par\smallskip\footnotesize $t$-statistics in parentheses. 1963-01 to 2026-06, $T = 762$.
All first-stage betas are re-estimated on this sample. Units and beta definitions as in
Table~\ref{tab:momfull}.
\end{table}

\begin{table}[H]\centering\small
\caption{Equation (3), full sample versus post-1963 $t$-statistics, 25 Size--momentum portfolios}
\begin{tabular}{lcc}
\toprule
 & Full sample & Post-1963 \\
\midrule
$\gamma_M$        & $-2.54$ & $-2.55$ \\
$\gamma_{size}$   & $-1.53$ & 0.19    \\
$\gamma_{ret212}$ & 2.77    & 4.13    \\
$\gamma_{SMB}$    & 0.56    & 0.68    \\
$\gamma_{UMD}$    & 0.35    & $-0.20$ \\
\bottomrule
\end{tabular}
\end{table}

\begin{tight}
  \item \textbf{Does the characteristics-vs.-covariances result change for momentum in the
  post-1963 sample?} The main conclusion does not change after 1963. In the combined
  equation (3), $\gamma_{ret212}$ remains strongly significant ($t = 4.13$), while
  $\gamma_{SMB}$ ($t = 0.68$) and $\gamma_{UMD}$ ($t = -0.20$) remain insignificant. The
  $\gamma_{ret212}$ coefficient changes only modestly from 0.00759 to 0.00804, while its
  FMB standard error falls from 0.00274 to 0.00195. The size evidence weakens: in equation
  (1), $\gamma_{size}$ falls from $-0.1319$ ($t = -4.35$) to $-0.0578$ ($t = -2.20$), and
  in equation (2) $\gamma_{SMB}$ loses significance after 1963 ($t = 3.45$ versus
  $t = 1.47$). In equation (3), the $\gamma_{size}$ $t$-statistic moves from $-1.53$ in
  the full sample to 0.19 post-1963, so the size effect largely disappears. Thus, the more
  robust result is the explanatory power of the past-return characteristic rather than
  size. As in the full sample, these results favor a characteristics interpretation
  relative to the measured SMB and UMD exposures, although they do not eliminate more
  complicated risk-based explanations.
  \item \textbf{Compare this to your findings for value in part (e).} Both value and
  momentum show fairly stable characteristics-versus-covariances conclusions across the
  full and post-1963 samples. For value, $\gamma_{B/M}$ in equation (3) is almost
  unchanged, moving from 0.4808 ($t = 4.39$) to 0.4876 ($t = 3.57$). $\gamma_{HML}$
  remains negative, while the absolute value of its $t$-statistic falls from 2.05 to 1.87.
  The size characteristic becomes essentially zero after 1963. For momentum,
  $\gamma_{ret212}$ also remains stable in magnitude, moving from 0.00759 ($t = 2.77$) to
  0.00804 ($t = 4.13$), while $\gamma_{UMD}$ remains insignificant in both samples. The
  post-1963 equation (3) UMD risk price ($-0.043$\% per month) is again far below UMD's
  average return over the same months (0.62\%).
  \item \textbf{Which premium appears more stable across sample periods, and what might
  that tell you?} It is difficult to conclude that one premium is clearly more stable
  overall. In equation (3), both characteristic coefficients are nearly unchanged across
  the two samples. In equation (1), however, BE/ME weakens more than ret212. The safest
  conclusion is that the characteristics-over-betas result survives the sample
  restriction for both value and momentum, while the size dimension is less stable.
  Because the post-1963 sample is contained within the full sample rather than being an
  independent out-of-sample period, this comparison should not be interpreted as a strong
  independent stability test.
\end{tight}

\paragraph{Statistical limitations.} This is a historical FMB horse race rather than a
real-time trading strategy. Full-period betas are generated regressors estimated using
future observations relative to early sample months. FMB standard errors use the standard
time-series SD divided by $\sqrt{T}$ and are not Shanken-corrected or HAC-adjusted.
Characteristics and estimated factor exposures are also highly correlated, so changes in
significance do not by themselves establish causality.

\end{document}
